\section{Discussion and Limitations}
\label{sec:discussion}

\subsection{Not relativism}

Declaring an observer does not mean that any story is acceptable. Facts,
causality, hashes, manifests, schemas, and verification results remain
load-bearing. The rule rejects an undeclared view from nowhere, not the
existence of facts.

\subsection{Not universal consensus}

Some systems require consensus, total ordering, or a central sequencer. The
observer-declared timeline rule does not forbid those designs. It says that
when a product presents a timeline to a human or agent, the view should still
declare the accepted facts and projection policy behind it.

\subsection{Adopter responsibility}

A schema can standardize vocabulary, but it cannot prove an adopter's runtime
correctness by itself. An adopter that claims a concrete timeline is complete,
fresh, or causally valid must provide product-specific evidence. This is why
KFD-4 should be assessed through KFD-2-style trust claims when used as a
release or product gate.

\subsection{Research status}

This is an initial design paper. The next version should replace design
intentions with concrete Kungfu implementation evidence: exported observer
perspectives, replay fixtures, failure cases, and measured recovery outcomes.
